\section{Related Work}
\subsection{Diffusion Models}
Diffusion models have been the de-facto choice in image and video generation tasks \cite{podell2023sdxl, gupta2024photorealistic, yang2024cogvideox}. 

Research in this field encompasses various aspects of improvement in     diffusion policy \cite{croitoru2023diffusion, liu2023flow}, model architecture \cite{bar2024lumiere,peebles2023scalable}, conditioning strategies \cite{zhang2023adding, khachatryan2023text2video}, among others.
For a comprehensive survey of image and video generation, we refer readers to~\cite{cao2024survey}. 
However, these methods are primarily designed for visual modality data. LiDAR data presents distinct challenges compared to visual data, as it describes a scene with unordered and unevenly distributed points. This fundamental difference makes it challenging to effectively leverage diffusion models for the LiDAR modality.



\subsection{LiDAR Scene Generation}

Recent LiDAR data generation approaches can be categorized as either unconditional or conditional generation. In the unconditional category, methods such as LiDARGen \cite{zyrianov2022learning}, R2DM \cite{nakashima2024lidar}, and LidarGRIT \cite{haghighi2024taming} train diffusion models directly on the pixel space, while RangeLDM additionally employs VAE \cite{kingma2013auto} to compress equirectangular images. In the conditional category, Text2LiDAR \cite{wu2024text2lidar} conditions the diffusion process on coarse textual scene descriptions. LiDM \cite{ran2024towards} investigates various conditioning inputs but applies each condition independently, limiting its comprehensive manipulation capabilities. Our object priors condition shares similarities with OLiDM \cite{yan2024olidm}, which utilizes synthetic objects as conditions for scene generation. However, OLiDM \cite{yan2024olidm} is designed to generate single LiDAR scenes, whereas our approach incorporates multiple multimodal conditions to generate \textit{sequential} LiDAR scenes—a non-trivial task requiring careful pipeline design. LidarDM \cite{zyrianov2024lidardm} decomposes the generation task into two stages, static and dynamic object synthesis. But this approach compromises the realism of object point clouds due to its simplistic modeling of 3D object structures. Moreover, 


We also see a growing research on jointly generating LiDAR point clouds and images. XDrive~\cite{xie2025xdrive} uses two diffusion models to generate image and LiDAR point clouds simultaneously, where latent LiDAR features and latent image features attend to each other via epipolar guidance. Albeit, XDrive \cite{xie2025xdrive} is limited to single LiDAR scene generation.
UniScene~\cite{li2024uniscene} indirectly enforces temporal consistency across sequential frames by mediating cross-modal interactions through implicit conditioning on 3D occupancy priors. Nevertheless, it exhibits limitations in synthesizing high-fidelity objects.


In summary, current methodologies in LiDAR generation field predominantly focus on 3D scene generation with coarse conditioning, neglecting the integration of complete temporal coherence with fine-grained controllability. To the best of our knowledge,  \sysname\ is the first unified framework to resolve the aforementioned limitations by integrating a equirectangular spatio-temporal diffusion model with multimodal conditioning capabilities.

